\par\noindent\SourceLeaf{57}%
\hspace*{\parindent}The other, however,

\[
\begin{aligned}
&n\cos\varphi\bigl(A^{(1)}\sin\varphi+2A^{(2)}\sin2\varphi+.\,.\,.\\&\quad+hA^{(h)}\sin h\varphi+(h+1)A^{(h+1)}\sin(h+1)\varphi+.\,.\,.\bigr),
\end{aligned}
\]

\noindent contains in each of its terms a product of $\cos\varphi$ and the sine of some multiple of $\varphi$. To combine it with the first into a single sum, we must use the known formula

\[
\cos\varphi\sin\lambda\varphi=\frac{1}{2}\sin(\lambda+1)\varphi+\frac{1}{2}\sin(\lambda-1)\varphi,
\]

\noindent to transform it. Thus the term in it containing $\sin\varphi$ arises from

\[
n\cos\varphi.2A^{(2)}\sin2\varphi
\]

\noindent and becomes $n.A^{(2)}\sin\varphi$. If, more generally, you look for terms in it containing $\sin(h+1)\varphi$, you must take the term

\[
hA^{(h)}\sin h\varphi.n\cos\varphi
\]

\noindent (and the first part of its expansion,

\[
\frac{nhA^{(h)}}{2}\sin(h+1)\varphi)
\]

\noindent together with the term

\[
(h+2)A^{(h+2)}\sin(h+2)\varphi.n\cos\varphi
\]

\noindent (and the latter part of its expansion,

\[
\frac{n(h+2)A^{(h+2)}}{2}\sin(h+1)\varphi);
\]

\Needspace{5\baselineskip}
\noindent After doing this, the coefficient of $\sin(h+1)\varphi$ in the series sought becomes

\[
\frac{nhA^{(h)}}{2}+\frac{n(h+2)A^{(h+2)}}{2}.
\]

\Needspace{7\baselineskip}
$\gamma$) Now since the series on one side ($\alpha$) equals the series taken together on the other side ($\beta$), the corresponding coefficients, that is, those multiplying the sine of the same angle here, must be equal. Thus on one side the coefficient of $\sin\varphi$ is

\[
=\nu nA-\nu\frac{nA^{(2)}}{2}
\]

\noindent and on the other it is $A^{(1)}+nA^{(2)}$. Equating them gives

\[
\nu nA-\nu\frac{nA^{(2)}}{2}=A^{(1)}+nA^{(2)},\qquad\mbox{whence}\qquad A^{(2)}=\frac{2(\nu nA-A^{(1)})}{(\nu+2)n}.
\]

